\section{Conclusion}
We asked whether a language model can reason directly over learned scientific features,
beyond describing them. On single objects and pairs it can, on both datasets. On questions
that require reasoning, the main result concerns training rather than architecture. A fully
fine-tuned narrator fails to find the odd star in a group, and fails more as the group grows.
A narrator trained with LoRA on a frozen base, on the same text and the same features,
succeeds in about two thirds of groups and does not lose accuracy with group size, and the
same comparison also favours LoRA on galaxies. Neither method was tuned, so a better recipe
for full fine-tuning may exist, but LoRA gives this result without one, and most of the gain
remains when the two share an optimizer schedule.

Latent tool use adds to this in a specific place. A small flow-matching model that receives
catalog labels as auxiliary data, and returns its output as a latent token, improves the
identification of outliers that never appeared as outliers in training ($0.62 \to 0.75$),
including when binary stars are excluded ($0.78 \to 0.88$). It adds little where the training
data offers precedent, and nothing on galaxies, where its output repeats what the narrator
already reads. Text is a poor substitute: models given the catalog numbers score below the
LoRA narrator, and frontier models given the raw observations zero-shot score lower still.
Verified templated text remains preferable to fluent LLM-written text as training data.

For a language model reading scientific data, these results suggest first adapting it with a
low-rank adapter, and then adding dedicated models with auxiliary data where the task goes
beyond what the training data covers. Why the frozen base preserves this ability, and how
frontier agents with code execution compare, remain open.

